Die Rekonstruktion zeitlich konsistenter 3D-Geometrie aus zwei bis vier Kameras
ist eine zentrale Herausforderung der Computer Vision mit direkter Relevanz für
Robotik und autonome Manipulation. Räumliche Foundation-Modelle erzielen hohe
geometrische Genauigkeit durch unabhängige Einzelbildrekonstruktion, erzeugen
jedoch hochfrequente zeitliche Schwankungen (Jitter) beim Zusammenführen zu einer
Sequenz. Bestehende Alternativen wie Geometrieverfeinerung und dedizierte
4D-Architekturen reduzieren Jitter auf Kosten geometrischer Genauigkeit, sodass
bislang kein Verfahren beides gleichzeitig erreicht.

Diese Arbeit verfolgt einen entkoppelten Ansatz: Ein leistungsfähiges räumliches
Modell rekonstruiert jedes Bild einzeln, eine leichtgewichtige zeitliche
Ausrichtungsstrategie überführt die Ergebnisse in ein gemeinsames
Koordinatensystem, und Confidence-Weighted Temporal Smoothing (Opt) unterdrückt
verbleibenden Jitter durch einen trainingsfreien Gaussfilter auf statischen
Szenenbereichen. Zeitliche Konsistenz und geometrische Genauigkeit stehen nur
dann im Widerspruch, wenn Konsistenzbedingungen auf dynamische Bereiche
ausgedehnt werden. Die Entkopplung beider Schritte vermeidet diesen Zielkonflikt
strukturell.

Sechs Rekonstruktionsmodelle werden auf DexYCB und
Hi4D in Zwei-, Drei- und Vier-Kamera-Konfigurationen evaluiert. Opt beseitigt
hochfrequenten Jitter vollständig und reduziert den mittleren Jitter um
60--91\,\% über vier Backbone-Modelle, während die Chamfer-Distanz bei drei der
vier Modelle unverändert bleibt oder sich verbessert. Die resultierende Pipeline
dominiert native 4D-Architekturen im Pareto-Sinne bei allen Kamerakonfigurationen.
Drift bleibt ein offenes Problem; eine global gekoppelte ankerbasierte Optimierung
ist der naheliegende nächste Schritt.